\documentclass[a4paper,12pt]{article}

\usepackage[utf8]{inputenc}
\usepackage[T1]{fontenc}
\usepackage[ngerman]{babel}
\usepackage{microtype}
\usepackage{geometry}
\usepackage{booktabs}
\usepackage{longtable}
\usepackage{enumitem}

\geometry{margin=2.5cm}

\title{Konnektoren im Deutschen: Arten, Funktion, Beispiele}
\date{}

\begin{document}
    \maketitle
    
    
    \section{Was sind Konnektoren?}
        Konnektoren sind W\"orter oder Wortgruppen, die S\"atze, Teils\"atze oder Satzglieder logisch miteinander verbinden.
        Sie zeigen Beziehungen wie Ursache, Folge, Gegensatz, Bedingung, Zeit, Zweck oder Ergänzung.
        
        \subsection{Wichtige Idee}
            Konnektoren sind \emph{keine} Kasus- oder Zeitform.
            Sie sind eine funktionale Kategorie: \emph{Verkn\"upfungselemente} im Satzbau.
    
    
    \section{Haupttypen von Konnektoren}
        Im Unterricht werden Konnektoren oft in drei große Gruppen gegliedert:
        \begin{itemize}[leftmargin=2em]
			\item koordinierende Konjunktionen (Verbposition bleibt meist gleich)
            \item subordinierende Konjunktionen / Subjunktionen (Verb am Ende im Nebensatz)
            \item adverbiale Konnektoren (meist Satzadverbien; Verb folgt direkt nach ihnen im Hauptsatz)
        \end{itemize}
    
    
    \section{Koordinierende Konjunktionen (Nebenordnung)}
        Diese Konjunktionen verbinden gleichrangige Einheiten (W\"orter, Satzteile, Haupts\"atze).
        Die Verbposition im zweiten Hauptsatz bleibt typischerweise wie im normalen Hauptsatz (Verb an Position 2).
        
        \subsection{Typische Konjunktionen}
            \begin{itemize}[leftmargin=2em]
				\item \textbf{und} (Addition)
                \item \textbf{oder} (Alternative)
                \item \textbf{aber} (Gegensatz)
                \item \textbf{denn} (Begründung)
                \item \textbf{sondern} (Korrektur nach Negation)
            \end{itemize}
        
        \subsection{Beispiele (erweitert)}
            \begin{itemize}[leftmargin=2em]
				\item \textbf{und:} Ich arbeite heute lange, \textbf{und} ich koche danach noch.
                \item \textbf{und:} Wir haben Kaffee gekauft \textbf{und} auch Tee.
                \item \textbf{oder:} Kommst du mit, \textbf{oder} bleibst du zu Hause?
                \item \textbf{oder:} Wir können um 10 Uhr starten \textbf{oder} um 11 Uhr.
                \item \textbf{aber:} Ich will trainieren, \textbf{aber} ich bin noch krank.
                \item \textbf{aber:} Das Problem ist schwierig, \textbf{aber} nicht unmöglich.
                \item \textbf{denn:} Ich gehe fr\"uher schlafen, \textbf{denn} ich bin m\"ude.
                \item \textbf{denn:} Wir bleiben hier, \textbf{denn} es regnet stark.
                \item \textbf{sondern:} Ich will nicht Pizza, \textbf{sondern} Pasta.
                \item \textbf{sondern:} Er ist nicht Student, \textbf{sondern} Forscher.
            \end{itemize}
        
        \subsection{Mini-Regel}
            \begin{itemize}[leftmargin=2em]
				\item Bei \textbf{und/oder/aber/denn/sondern} steht im zweiten Hauptsatz das Verb in der Regel wieder an Position 2:
                \item Beispiel: Ich bleibe zu Hause, aber \textbf{ich gehe} sp\"ater einkaufen.
            \end{itemize}
    
    
    \section{Subordinierende Konjunktionen / Subjunktionen (Unterordnung)}
        Diese Konnektoren leiten einen Nebensatz ein.
        Im Nebensatz steht das konjugierte Verb typischerweise \textbf{am Ende}.
        
        \subsection{Sehr h\"aufige Subjunktionen}
            \begin{longtable}{@{}ll@{}}
                \toprule
                Konnektor & Typische Beziehung \\
                \midrule
                \endhead
                \textbf{weil} & Ursache (Begründung) \\
                \textbf{obwohl} & Gegensatz (Konzession) \\
                \textbf{wenn} & Bedingung / Zeit (je nach Kontext) \\
                \textbf{dass} & Inhaltssatz (was jemand denkt/sagt) \\
                \textbf{damit} & Zweck (Ziel) \\
                \textbf{w\"ahrend} & Zeit / Gegensatz (je nach Kontext) \\
                \textbf{bevor} & Zeit (vorher) \\
                \textbf{nachdem} & Zeit (nachher) \\
                \bottomrule
            \end{longtable}
        
        \subsection{Beispiele (erweitert)}
            
            \subsubsection{weil (Ursache)}
            \begin{itemize}[leftmargin=2em]
					\item Ich bleibe heute zu Hause, \textbf{weil ich krank bin}.
                    \item Wir nehmen ein Taxi, \textbf{weil der Bus nicht mehr f\"ahrt}.
                    \item Ich trinke Wasser, \textbf{weil ich Durst habe}.
                \end{itemize}
            
            \subsubsection{obwohl (Gegensatz)}
            \begin{itemize}[leftmargin=2em]
					\item Ich gehe spazieren, \textbf{obwohl es regnet}.
                    \item Er hat die Aufgabe geschafft, \textbf{obwohl sie schwer war}.
                    \item Wir treffen uns, \textbf{obwohl wir wenig Zeit haben}.
                \end{itemize}
            
            \subsubsection{wenn (Bedingung oder Zeit)}
            \begin{itemize}[leftmargin=2em]
					\item \textbf{Bedingung:} Ich helfe dir, \textbf{wenn du mich brauchst}.
                    \item \textbf{Bedingung:} Wenn du willst, \textbf{k\"onnen wir} morgen telefonieren. (Nebensatz zuerst)
                    \item \textbf{Zeit (immer wenn):} Ich werde nerv\"os, \textbf{wenn ich pr\"ufe}.
                \end{itemize}
            
            \subsubsection{dass (Inhaltssatz)}
            \begin{itemize}[leftmargin=2em]
					\item Ich glaube, \textbf{dass er heute keine Zeit hat}.
                    \item Sie sagt, \textbf{dass die Ergebnisse morgen kommen}.
                    \item Wir hoffen, \textbf{dass alles gut l\"auft}.
                \end{itemize}
            
            \subsubsection{damit (Zweck)}
            \begin{itemize}[leftmargin=2em]
					\item Ich lerne jeden Tag, \textbf{damit ich die Pr\"ufung bestehe}.
                    \item Wir schreiben alles auf, \textbf{damit wir nichts vergessen}.
                    \item Er spricht langsam, \textbf{damit alle ihn verstehen}.
                \end{itemize}
            
            \subsubsection{w\"ahrend (Zeit oder Kontrast)}
            \begin{itemize}[leftmargin=2em]
					\item \textbf{Zeit:} Ich h\"ore Musik, \textbf{w\"ahrend ich arbeite}.
                    \item \textbf{Kontrast:} W\"ahrend er gern fr\"uh aufsteht, \textbf{schlafe ich} lieber lang. (Nebensatz zuerst)
                \end{itemize}
            
            \subsubsection{bevor / nachdem (Zeit)}
            \begin{itemize}[leftmargin=2em]
					\item Ich trinke Kaffee, \textbf{bevor ich anfange zu arbeiten}.
                    \item Bevor wir starten, \textbf{pr\"ufen wir} die Batterien.
                    \item Nachdem ich gegessen habe, \textbf{mache ich} einen kurzen Spaziergang.
                    \item Ich rufe dich an, \textbf{nachdem ich fertig bin}.
                \end{itemize}
        
        \subsection{Mini-Regel: Verb am Ende}
            \begin{itemize}[leftmargin=2em]
				\item \textbf{weil/obwohl/wenn/dass/damit/...} \(\rightarrow\) Nebensatz: Verb \textbf{am Ende}
                \item Beispiel: ..., weil ich heute \textbf{keine Zeit habe}.
            \end{itemize}
    
    
    \section{Adverbiale Konnektoren (Satzadverbien / Konnektoradverbien)}
        Diese Konnektoren verbinden meist ganze S\"atze und stehen oft am Satzanfang.
        Wichtig: Wenn sie am Anfang stehen, kommt das konjugierte Verb im Hauptsatz direkt danach (Verb Position 2).
        
        \subsection{H\"aufige adverbiale Konnektoren}
            \begin{longtable}{@{}ll@{}}
                \toprule
                Konnektor & Beziehung \\
                \midrule
                \endhead
                \textbf{deshalb / deswegen / daher} & Folge (aus einem Grund) \\
                \textbf{trotzdem} & Gegensatz (trotz X) \\
                \textbf{au\ss erdem} & Zusatz / Ergänzung \\
                \textbf{dann} & zeitliche Folge \\
                \textbf{allerdings} & Einschr\"ankung \\
                \textbf{zum Beispiel} & Beispiel \\
                \bottomrule
            \end{longtable}
        
        \subsection{Beispiele (erweitert)}
            
            \subsubsection{deshalb / deswegen / daher (Folge)}
            \begin{itemize}[leftmargin=2em]
					\item Es regnet stark. \textbf{Deshalb bleiben} wir zu Hause.
                    \item Ich war krank. \textbf{Deswegen habe} ich nicht trainiert.
                    \item Der Bus kommt nicht. \textbf{Daher nehmen} wir ein Taxi.
                \end{itemize}
            
            \subsubsection{trotzdem (Gegensatz)}
            \begin{itemize}[leftmargin=2em]
					\item Es ist kalt. \textbf{Trotzdem gehe} ich raus.
                    \item Ich habe wenig Zeit. \textbf{Trotzdem schreibe} ich noch eine Seite.
                    \item Die Aufgabe ist schwer. \textbf{Trotzdem mache} ich weiter.
                \end{itemize}
            
            \subsubsection{au\ss erdem (Zusatz)}
            \begin{itemize}[leftmargin=2em]
					\item Ich bin m\"ude. \textbf{Au\ss erdem muss} ich fr\"uh aufstehen.
                    \item Wir haben Brot gekauft. \textbf{Au\ss erdem brauchen} wir noch Milch.
                    \item Er ist kompetent. \textbf{Au\ss erdem arbeitet} er sehr genau.
                \end{itemize}
            
            \subsubsection{dann (zeitliche Folge)}
            \begin{itemize}[leftmargin=2em]
					\item Ich arbeite zuerst zwei Stunden, \textbf{dann mache} ich eine Pause.
                    \item Wir essen, \textbf{dann gehen} wir spazieren.
                    \item Ich rufe dich an; \textbf{dann k\"onnen} wir alles kl\"aren.
                \end{itemize}
            
            \subsubsection{allerdings (Einschr\"ankung)}
            \begin{itemize}[leftmargin=2em]
					\item Das klingt gut. \textbf{Allerdings ist} es teuer.
                    \item Ich w\"urde gern kommen. \textbf{Allerdings habe} ich schon einen Termin.
                    \item Der Plan ist sinnvoll. \textbf{Allerdings brauchen} wir mehr Daten.
                \end{itemize}
        
        \subsection{Mini-Regel: Verb direkt nach dem Konnektor}
            \begin{itemize}[leftmargin=2em]
				\item \textbf{Deshalb/Trotzdem/Au\ss erdem/...} am Satzanfang \(\rightarrow\) Verb folgt sofort:
                \item Beispiel: Deshalb \textbf{gehe} ich jetzt.
            \end{itemize}
    
    
    \section{Schneller Vergleich (Merkschema)}
        \begin{longtable}{@{}p{0.28\textwidth}p{0.34\textwidth}p{0.30\textwidth}@{}}
            \toprule
            Typ & Beispiel-Konnektoren & Satzbau-Kernregel \\
            \midrule
            \endhead
            Koordinierende Konjunktionen &
            und, oder, aber, denn, sondern &
            2. Hauptsatz: Verb bleibt typischerweise Position 2 \\
            \addlinespace
            Subjunktionen (Nebensatz) &
            weil, obwohl, wenn, dass, damit, bevor, nachdem &
            Nebensatz: konjugiertes Verb am Ende \\
            \addlinespace
            Adverbiale Konnektoren &
            deshalb, trotzdem, au\ss erdem, dann, allerdings &
            Wenn am Anfang: Verb direkt danach (Position 2) \\
            \bottomrule
        \end{longtable}

\end{document}